\chapter{Mathematical Foundations and Proofs}

\section{Purpose and Standing Assumptions}

The main text repeatedly presents a result informally and names, honestly, what remains to be shown before the informal statement becomes a theorem. This appendix takes up that debt directly. Every result below is conditional: each is stated as holding whenever certain hypotheses about 𝓜, C, and H are satisfied, hypotheses the main text motivated but did not verify for any particular implementation. Stating results this way is not a retreat from rigor. It is what rigor actually requires once a framework is meant to cover many different applications, a courtyard, a robot, a civilization, whose configuration spaces will not all share the same fine structure. The standing assumptions, used without restatement throughout this appendix unless explicitly relaxed, are as follows. 𝓜 is a complete metric space under the metric g of Chapter 18.3. C ⊆ 𝓜 is the constraint manifold of Chapter 18.5. H, the energy functional of Chapter 19.2, is bounded below.

\section{Existence of Constraint Projections}

Theorem B.1 (Existence). If C is closed in 𝓜, then for every Y ∈ 𝓜 there exists a point \(\Pi_C\)(Y) ∈ C minimizing the distance to Y, that is, satisfying ‖\(\Pi_C\)(Y) − Y‖ = \(\inf_{Z \in C}\) ‖Z − Y‖.

Proof sketch. Let d = \(\inf_{Z \in C}\) ‖Z − Y‖ and choose a sequence (\(Z_n\)) in C with ‖\(Z_n\) − Y‖ → d, a minimizing sequence. Because 𝓜 is complete, it suffices to show (\(Z_n\)) has a convergent subsequence whose limit lies in C; the limit then achieves the infimum by continuity of the distance function. If 𝓜 is additionally locally compact near Y, or if the sublevel set \{Z ∈ C : ‖Z − Y‖ ≤ d + 1\} is compact, a subsequence converges by the Bolzano–Weierstrass property, and closedness of C guarantees the limit is itself a member of C rather than merely a boundary point outside it. This is the direct method of the calculus of variations, applied here to a distance functional rather than to H directly; Chapter 19.6 already named the two conditions this proof actually needs, coercivity, supplying the compactness of the relevant sublevel set, and closedness of C, guaranteeing the limit remains admissible.

This theorem is conditional on C being closed and on some compactness condition holding near the minimizing sequence. Neither is guaranteed by this book's definition of C as an intersection \(C_S\) ∩ \(C_z\) ∩ \(C_R\) ∩ \(C_A\) ∩ ⋯; closedness of an intersection of closed sets is automatic, but whether each individual member is closed, and whether the relevant compactness holds, depends on details, the precise topology of \(L_t(x)\)'s space of live possibilities, the precise structure of residue accumulation, that this book has motivated without fully specifying. Section B.9 returns to this.

\section{On the Uniqueness of Projections}

It is worth being direct about a claim this appendix does not make. Uniqueness of \(\Pi_C\)(Y) is not guaranteed by Theorem B.1 alone; it additionally requires something like convexity of C, or strict convexity of the norm defining g, conditions this book has never asserted and has no particular reason to expect. Where C is not convex, and a constraint manifold built from several independently motivated conditions, entropy, appearance consistency, residue consistency, affordance-realizability, has no general reason to be, two or more equidistant points of C can exist for a single proposal, each an equally valid nearest admissible state. This is not a defect in the theorem. It is a structural feature the main text already built machinery to handle: Chapter 15's account of refusal and arbitration among simultaneous proposals is, read carefully, partly a response to exactly this possibility, that projection alone may not determine a unique outcome and something further, commit-time arbitration, is needed to select among admissible candidates when more than one exists.

\section{Convergence of the Write Cycle}

Write Ψ for the combined map taking a committed state to its successor under one full cycle, Ψ(M) = \(\Pi_C\)(M + \(\Delta\psi\)(M)), treating the proposal \(\Delta\psi\) as itself a (possibly stochastic) function of the current state for the purposes of this section.

Theorem B.2 (Conditional convergence). If Ψ is a contraction on C, that is, if there exists k < 1 such that ‖Ψ(M) − Ψ(M′)‖ ≤ k‖M − M′‖ for all M, M′ ∈ C, then Ψ has a unique fixed point \(M^*\) ∈ C, and the sequence M, Ψ(M), Ψ(Ψ(M)), … converges to \(M^*\) from any starting point.

Proof sketch. This is the Banach fixed-point theorem, applied to Ψ on the complete metric space C (closed subsets of complete spaces are themselves complete). The standard proof shows the sequence of iterates is Cauchy, using the contraction property to bound ‖\(\Psi^n\)(M) − \(\Psi^{n+1}\)(M)‖ by \(k^n\) times a constant, and completeness supplies a limit; contractivity then shows the limit is fixed and unique.

This theorem is considerably stronger than anything the main text actually established. Chapter 16 argued that fixed points exist as a matter of definition, wherever a trajectory happens to settle, without claiming every trajectory settles, or that settlement is unique when it occurs. The contraction hypothesis of Theorem B.2 is a genuine further assumption, plausible for a system whose intrinsic dynamics \(F_{\mathrm{phys}}\) are strongly dissipative in Chapter 19's sense but not established for this book's \(\Delta\psi\) in general, particularly once \(F_\psi\)'s learned, expressive contribution is included. Where Ψ is not a contraction, weaker conclusions are available under weaker hypotheses, for instance that Ψ is merely non-expansive (k ≤ 1) and C is compact, which guarantees existence of a fixed point by a Schauder-type argument without guaranteeing uniqueness or convergence from an arbitrary start. This book's own examples, the fountain especially, suggest the compact-but-not-contracting case is closer to the typical one: a dynamic fixed point in Chapter 16 and Chapter 20's sense is not approached and settled into once, the way a contraction's fixed point is, but continuously revisited by an orbit that never stops moving.

\section{Stability of Fixed Points}

Fix a fixed point \(M^*\) and suppose C is, near \(M^*\), a smooth submanifold of 𝓜, so that the tangent space \(T_{M^*}C\) is well defined.

Proposition B.3 (Second-order stability test). Let \(\operatorname{Hess}H|_C\) denote the Hessian of H restricted to C at \(M^*\). If \(\operatorname{Hess}H|_C\) is positive definite on \(T_{M^*}C\), then \(M^*\) is a stable fixed point in Chapter 16.5's sense: sufficiently small admissible perturbations of \(M^*\) are followed, under the write cycle, by a return toward \(M^*\). If \(\operatorname{Hess}H|_C\) has a negative eigenvalue on \(T_{M^*}C\), then \(M^*\) is unstable: some admissible perturbation, however small, is followed by movement away from \(M^*\).

Proof sketch. This is the standard second-derivative test for a critical point of a constrained optimization problem, transported to this book's dynamical setting. Because \(M^*\) is a fixed point, it is a critical point of H restricted to C (Chapter 19.5); the sign of the Hessian at a critical point determines, to second order, whether nearby points have higher or lower energy. Since the write cycle's dissipative component moves the field toward lower energy within C (Chapter 19.3), a direction of strictly higher energy is a direction the dynamics push back from, giving stability, while a direction of lower energy is a direction the dynamics continue to move along, giving instability. The argument is local and says nothing about behavior far from \(M^*\); a fixed point can be locally stable under this test and still lie within a region where global dynamics behave in ways this proposition does not address.

\section{Sheaf Gluing}

Proposition B.4 (Gluing criterion). Let \(F_{\mathrm{field}}\) be the field-coherence sheaf of Chapter 21.2 over a cover 𝒢 of X. A compatible family of local sections, one admissible configuration per patch, agreeing on all pairwise overlaps, glues to a unique global section if and only if the corresponding first Čech cohomology group vanishes, \(H^1\)(𝒢, \(F_{\mathrm{field}}\)) = 0.

This is a standard fact of sheaf theory rather than a result original to this book; it is included here because the main text invokes it, in Chapter 21.3, without stating it in this form. The forward direction, vanishing cohomology implies gluing, follows from the long exact sequence relating local sections, global sections, and the obstruction cocycles Čech theory constructs from pairwise overlaps; a nontrivial class in \(H^1\) is, by construction, exactly a compatible-on-overlaps family that cannot be extended to a global section, since resolving it would require the coboundary of some cochain that the nontrivial class witnesses does not exist. This gives the main text's claim that hallucination, in this book's spatial sense, is precisely a nonvanishing \(H^1\), a genuine consequence of standard cohomological gluing theory rather than an analogy borrowed from it.

\section{Conservation from Symmetry}

Theorem B.5 (Noether's theorem, as used in Chapter 22). Let L be a Lagrangian on the field's configuration space, and suppose a one-parameter family of transformations \(M_s\), with \(M_0\) = M, leaves L invariant, L(\(M_s\), \(\dot{M}_s\)) = L(M, \(\dot{M}\)) for all s. Then the quantity

Q = \(\textstyle\sum_i\) (\(\partial L/\partial \dot{M}_i\)) (\(dM_{s,i}/ds\))\(\big|_{s=0}\)

is constant along any trajectory satisfying the Euler–Lagrange equations derived from L.

Proof sketch. Differentiate the invariance condition with respect to s at s = 0 to obtain a relation among \(\partial L/\partial M\), \(\partial L/\partial \dot{M}\), and the generator of the symmetry, d\(M_s\)/ds. Substituting the Euler–Lagrange equations, which express \(\partial L/\partial M\) in terms of the time derivative of \(\partial L/\partial \dot{M}\), converts this relation into dQ/dt = 0. This is the standard proof, unmodified by anything specific to semantic fields; what is specific to this book is only the identification of the three symmetries in Chapter 22.3, time-translation, domain-relabeling of commutative commits, and domain-translation, as genuine invariances of the Lagrangian L introduced in Chapter 19.2, an identification this appendix does not independently re-verify but that the main text's construction of L was explicitly built to support.

\section{Commutativity of Independent Updates}

Proposition B.6. Let \(\Delta\psi_a\) and \(\Delta\psi_b\) be two proposals with disjoint support, meaning each is identically zero outside a region of X that does not intersect the other's. Suppose further that C is a product constraint, in the sense that admissibility on disjoint regions can be checked independently, and that \(\Pi_C\) acts componentwise on such disjoint regions. Then

\(\Pi_C\)(\(\Pi_C\)(M + \(\Delta\psi_a\)) + \(\Delta\psi_b\)) = \(\Pi_C\)(\(\Pi_C\)(M + \(\Delta\psi_b\)) + \(\Delta\psi_a\)),

that is, committing the two proposals in either order yields the same result.

Proof sketch. Because the two proposals have disjoint support and the write operator is additive (Chapter 12.3), M + \(\Delta\psi_a\) + \(\Delta\psi_b\) is unambiguous regardless of the order the two terms are added in. Because C is assumed to be a product constraint and \(\Pi_C\) acts componentwise on disjoint regions, projecting this sum is equivalent to projecting each region against its own admissibility condition independently, and this per-region projection does not depend on whether the other region's proposal has already been incorporated, since by hypothesis the two regions do not interact. This gives the result claimed in Chapter 15.3 without further argument, provided its two hypotheses, disjoint support and a product constraint structure, actually hold. The second hypothesis is the substantive one: several members of C, most notably \(C_A\), affordance-realizability, and any global conservation law from Chapter 22, are not obviously product constraints in this sense, since a conservation law is precisely a condition relating what happens in one region to what happens elsewhere. Where a member of C fails to be a product constraint, Proposition B.6 does not apply as stated, and commutativity of disjoint updates would need to be re-examined against that member specifically.

\section{What Remains Open}

Every result in this appendix is conditional, and it is worth collecting, in one place, exactly which conditions have not been verified for this book's own constructions. Closedness of C (B.2) depends on the topology of each of its members, established informally in Chapters 8, 9, 11, and 21 but never given a topology precise enough to check closedness against. Coercivity of H (B.2, B.4) has been assumed, not derived, from the specific Lagrangian of Chapter 19.2. The contraction or non-expansiveness of the full write cycle Ψ (B.4) is plausible for the dissipative component \(F_{\mathrm{phys}}\) and unexamined for the learned component \(F_\psi\). Whether C's individual members are product constraints in the sense B.8 requires is, for at least two of them, affordance-realizability and any conservation law, doubtful as stated and would need a genuinely different argument, likely one that bounds rather than eliminates the interaction between disjoint regions. None of this undermines the results as conditional theorems; each is a correct statement of what follows from its hypotheses. It does mean that a reader implementing this architecture for a specific application inherits an obligation this appendix has been explicit about rather than hidden: verifying, for that application's specific 𝓜, C, and H, that the hypotheses these theorems require actually hold, before relying on the guarantees they would otherwise supply.
